% !TeX root = ../colloquium-2.tex

\phantomsection\label{sec_Proofs02}
\begin{proof}[Доказательство 1-й теоремы Вейерштрасса. Доказательство 2-й теоремы Вейерштрасса]
    \subsection*{Доказательство 1-й теоремы Вейерштрасса.}
    \begin{note}[amber]
        Пусть функция $f$ определена и непрерывна на отрезке $[a, b]$. Тогда она ограничена на этом отрезке. \\
        \begin{note}[purple]
            Доказательство.Если функция $f$ не ограничена на отрезке, то для каждого натурального n существует такая точка $x_n \in [a,b] $,
            что $|f(x_n)| > n$. Все элементы последовательности $\{x_n\}^{+\infty}_{n = 1}$ принадлежат отрезку
            $[a, b]$, то есть последовательность ограничена, поэтому по теореме Больцано – Вейерштрасса она имеет хотя бы одну cходящуюся к некоторому числу c
            подпоследовательность
            $c_k = a_{n_k}$ , причём $c \in[a, b]$ в силу предельного перехода в неравенствах. При этом
            $|f(c_k)| > n_k$, поэтому в точке c у функции $f$ не существует конечного предела, что противоречит непрерывности этой функции.
        \end{note}
    \end{note}

    \subsection*{Доказательство 2-й теоремы Вейерштрасса.}
    \begin{note}[amber]
        Пусть функция $f$ определена и непрерывна на отрезке $[a, b]$. Тогда существуют такие точки  $x_1, x_2 \in [a,b]$, что
        \[
        f(x_1)=m=\inf_{a\le x\le b} f(x), \ f(x_2)=M=\sup_{a\le x\le b} f(x).
        \]\\
        \begin{note}[purple]
            Доказательство.  Существование точных верхней и нижней граней функции $f$ следует из того, что по первой теореме Вейерштрасса функция $f$ ограничена на отрезке $[a,b]$, то есть множество её значений ограничено, а тогда существует 
            точная верхняя и точная нижняя грани её множества значений.
            Доказательство проведём для супремума. Предположим, что точка $x_2$ не найдётся. Тогда
            \[
            M-f(x)>0 \qquad \forall x\in [a,b].
            \]
            Рассмотрим функцию $g$, которая в точке $x\in [a,b]$ принимает значение
            \[
            g(x)=\frac{1}{M-f(x)}.
            \]
            В силу локальных свойств непрерывных функций $g$ непрерывна в каждой точке отрезка $[a,b]$, а тогда она непрерывна на всём отрезке по определению, что даёт её ограниченность на всём отрезке по 1-й теореме Вейерштрасса. С другой стороны, по определению супремума
            \[
            \forall \varepsilon>0 \ \exists x_\varepsilon\in [a,b] :
            f(x_\varepsilon)>M - \varepsilon.
            \]
            При $\varepsilon=\frac{1}{n}$ получим, что $\forall n\in\mathbb{N}$
            \[
            g(x_{1/n})>n,
            \]
            то есть функция $g$ не ограничена на отрезке. Таким образом, пришли к противоречию.
        \end{note}
    \end{note}
\end{proof}
